\documentclass{report}

\input{../LatexTemp/preamble.tex}
\input{../LatexTemp/macros}
\input{../LatexTemp/letterfonts}

\usepackage[utf8]{inputenc}

\setlength{\parindent}{0pt}

\title{Esercizio: Ludoteca — Correttezza della riduzione $k$-VCOL $\leq_P$ Ludoteca}
\author{Giovanni Palma}
\date{}

\begin{document}

\maketitle

\section*{Costruzione della riduzione}

Sia $\langle G = (V, E), k\rangle$ un'istanza di $k$-VCOL. Definiamo $f(\langle G, k\rangle) = \langle B, S, G_{\text{lud}}, p\rangle$:
\begin{itemize}
  \item $B = V$ (un bambino per nodo)
  \item $S = \{s_1, \ldots, s_k\}$ ($k$ stanze, una per colore)
  \item $G_{\text{lud}} = \{g_e : e \in E\}$ (un gioco per arco)
  \item $p(v) = \{g_e : e \in E,\; v \in e\}$ (il bambino $v$ desidera i giochi corrispondenti agli archi incidenti)
\end{itemize}

\textbf{Proprietà di costruzione (fondamentale).} Per ogni $u, v \in V$:
\[
  p(u) \cap p(v) \neq \emptyset \iff \{u, v\} \in E.
\]
Infatti $g_e \in p(u) \cap p(v)$ sse $u, v \in e$, cioè $e = \{u, v\}$. Quindi \emph{conflitto in Ludoteca} $\iff$ \emph{adiacenza in $G$}.

\section*{Correttezza}

Mostriamo $\langle G, k\rangle \in k\text{-VCOL} \iff f(\langle G, k\rangle) \in \text{Ludoteca}$.

\subsection*{($\Rightarrow$) $G$ $k$-colorabile $\Rightarrow$ istanza Ludoteca "sì"}

Sia $c : V \to \{1, \ldots, k\}$ una $k$-colorazione valida di $G$. Definiamo l'assegnazione
\[
  a : B \to S, \qquad a(v) := s_{c(v)}.
\]
Mostriamo che $a$ è senza conflitti. Siano $u, v \in B$ con $a(u) = a(v)$:
\[
  a(u) = a(v) \;\Rightarrow\; s_{c(u)} = s_{c(v)} \;\Rightarrow\; c(u) = c(v) \;\overset{c \text{ valida}}{\Longrightarrow}\; \{u, v\} \notin E.
\]
Per la proprietà di costruzione, $\{u, v\} \notin E \Rightarrow p(u) \cap p(v) = \emptyset$. Nessun conflitto, quindi $f(\langle G, k\rangle) \in \text{Ludoteca}$. \qed

\subsection*{($\Leftarrow$) istanza Ludoteca "sì" $\Rightarrow$ $G$ $k$-colorabile}

Sia $a : B \to S$ assegnazione senza conflitti. Definiamo la colorazione
\[
  c : V \to \{1, \ldots, k\}, \qquad c(v) := i \quad \text{sse} \quad a(v) = s_i.
\]
Mostriamo che $c$ è una $k$-colorazione valida. Sia $\{u, v\} \in E$:
\[
  \{u, v\} \in E \;\Rightarrow\; g_{\{u,v\}} \in p(u) \cap p(v) \;\Rightarrow\; u, v \text{ in conflitto}.
\]
Poiché $a$ è senza conflitti, $a(u) \neq a(v)$, da cui $c(u) \neq c(v)$. Quindi $c$ è $k$-colorazione valida, e $\langle G, k\rangle \in k\text{-VCOL}$. \qed

\section*{Calcolabilità polinomiale}

$f$ scandisce $E$ (per creare i giochi) e per ogni $v \in V$ enumera gli archi incidenti (per costruire $p(v)$). Costo $O(|V| + |E|)$, polinomiale in $\|\langle G, k\rangle\|$.

\section*{Conclusione}

Ludoteca $\in NP$ (certificato = assegnazione $a$, verifica in $O(|B|^2 \cdot |G_{\text{lud}}|)$). $3$-VCOL è NP-completo, e abbiamo dimostrato $3\text{-VCOL} \leq_P \text{Ludoteca}$ (caso $k = 3$). Quindi \textbf{Ludoteca è NP-completo}.

\end{document}
